% =========================================================
% Dense Subsets and Separability
% Source: volume-iii/analysis/metric-spaces/notes/metric-spaces/
%
% Dependencies:
%   notes-metric-space.tex       — def:metric, def:open-ball
%   notes-limit-points.tex       — def:closure
%   notes-completeness.tex       — completion idea
%   notes-real-analysis-bridge.tex — P([a,b]) / C([a,b]) reference
% =========================================================

% ---------------------------------------------------------
\subsubsection{Dense Subsets and Separability}
\label{sec:dense-separable}
% ---------------------------------------------------------

% ---- Toolkit ----
\begin{tcolorbox}[colback=gray!6, colframe=gray!40, arc=2pt,
  left=6pt, right=6pt, top=4pt, bottom=4pt,
  title={\small\textbf{Dense Subsets and Separability --- Quick Reference}},
  fonttitle=\small\bfseries]
\small
\begin{tabular}{l l l l}
\toprule
\textbf{Label} & \textbf{Statement} & \textbf{Proof method} & \textbf{Detail} \\
\midrule
D\ref{def:dense}
  & $D$ dense in $X$ iff $\overline{D} = X$
  & Definition
  & \hyperref[def:dense]{$\downarrow$} \\[4pt]
D\ref{def:separable}
  & $(X,d)$ separable iff $\exists$ countable dense $D \subseteq X$
  & Definition
  & \hyperref[def:separable]{$\downarrow$} \\[4pt]
P\ref{prop:dense-equiv}
  & $\overline{D}=X \iff \forall x \in X,\, \forall r>0,\, B_r(x)\cap D \neq \varnothing$
  & Closure characterization
  & \hyperref[prop:dense-equiv]{$\downarrow$} \\[4pt]
P\ref{prop:Q-dense-R}
  & $\mathbb{Q}$ is dense in $\mathbb{R}$; $\mathbb{R}$ is separable
  & Archimedean property
  & \hyperref[prop:Q-dense-R]{$\downarrow$} \\[4pt]
P\ref{prop:Rn-separable}
  & $\mathbb{R}^n$ is separable (dense subset $\mathbb{Q}^n$)
  & Component-wise
  & \hyperref[prop:Rn-separable]{$\downarrow$} \\[4pt]
P\ref{prop:discrete-not-sep}
  & Uncountable discrete space is not separable
  & Counterexample
  & \hyperref[prop:discrete-not-sep]{$\downarrow$} \\[4pt]
\bottomrule
\end{tabular}
\end{tcolorbox}

% ---------------------------------------------------------
\paragraph{Motivation.}
When we said that the polynomials $\mathbb{P}([a,b])$ are ``completed'' to the
continuous functions $C([a,b])$, we were using an idea that deserves to be
stated explicitly: every continuous function can be \emph{approximated
arbitrarily closely} by polynomials.
That approximation statement is exactly what \emph{density} means.

More broadly, density answers the question: can every point in a space be
reached as a limit of points from some simpler, smaller subset?
If yes, then to understand the whole space it can be enough to understand
the dense subset --- convergence, limits, and completions all depend on
this idea.

% ---- Definition: Dense ----
\begin{tcolorbox}[colback=propbox, colframe=propborder, arc=2pt,
  left=6pt, right=6pt, top=4pt, bottom=4pt,
  title={\small\textbf{Definition (Dense Subset)}},
  fonttitle=\small\bfseries]
\begin{definition}[Dense Subset]
\label{def:dense}
Let $(X,d)$ be a metric space and $D \subseteq X$.
We say $D$ is \emph{dense in $X$} if
\[
\overline{D} = X.
\]
\end{definition}
\end{tcolorbox}

\begin{remark*}[Fully quantified form]
\[
D \text{ is dense in } X
\;\iff\;
\forall\, x \in X,\quad x \in \overline{D}.
\]
That is, every point of $X$ is either in $D$ or is a limit point of $D$.
\end{remark*}

\begin{remark*}[Intuition]
Think of $D$ as a set of ``approximations.''
Saying $D$ is dense in $X$ means: no matter where you stand in $X$,
you are either already in $D$, or you can be approached arbitrarily closely
by points of $D$.
There are no ``gaps'' in $X$ that $D$ cannot reach.

The classic example is $\mathbb{Q}$ dense in $\mathbb{R}$: every real
number, rational or irrational, can be approximated to any precision by
rationals.
\end{remark*}

% ---- Proposition: equivalent characterization ----

\begin{proposition}[Geometric characterization of density]
\label{prop:dense-equiv}
Let $(X,d)$ be a metric space and $D \subseteq X$. The following are equivalent:
\begin{enumerate}[(i)]
  \item $D$ is dense in $X$ (i.e.\ $\overline{D} = X$).
  \item Every non-empty open ball in $X$ meets $D$:
        $\forall\, x \in X,\; \forall\, r > 0,\; B_r(x) \cap D \neq \varnothing$.
  \item Every non-empty open set $U \subseteq X$ meets $D$:
        $\forall\, U \subseteq X$ open and non-empty, $U \cap D \neq \varnothing$.
\end{enumerate}
\end{proposition}

\begin{proof}
(i) $\Rightarrow$ (ii): Let $x \in X$ and $r > 0$.
Since $\overline{D} = X$, we have $x \in \overline{D}$,
so by the geometric characterization of closure,
every ball $B_r(x)$ meets $D$.

(ii) $\Rightarrow$ (iii): Every non-empty open set $U$ contains some ball
$B_r(x) \subseteq U$, which by (ii) meets $D$.
Hence $U \cap D \supseteq B_r(x) \cap D \neq \varnothing$.

(iii) $\Rightarrow$ (i): Let $x \in X$.
For any $r > 0$, the ball $B_r(x)$ is open and non-empty,
so by (iii) it meets $D$.
Hence $x \in \overline{D}$ by the geometric characterization of closure.
Since $x$ was arbitrary, $\overline{D} = X$.
\end{proof}

\begin{remark*}[Why (iii) is the most useful form]
Condition (iii) --- that $D$ meets every open set --- is the form
most often used in proofs.
It frees you from having to work with specific balls;
any open neighborhood of any point in $X$ will do.
\end{remark*}

% ---- Examples ----

\begin{proposition}[$\mathbb{Q}$ is dense in $\mathbb{R}$]
\label{prop:Q-dense-R}
The rationals $\mathbb{Q}$ are dense in $(\mathbb{R}, |\cdot|)$.
Equivalently, $\mathbb{R}$ is the closure of $\mathbb{Q}$.
\end{proposition}

\begin{proof}
Let $x \in \mathbb{R}$ and $r > 0$.
By the Archimedean property, there exists a rational $q$ with
$|q - x| < r$, i.e.\ $q \in B_r(x)$.
So every open ball in $\mathbb{R}$ contains a rational, and
$\overline{\mathbb{Q}} = \mathbb{R}$.
\end{proof}

\begin{remark*}[Consequence: the irrationals are also dense]
The irrationals $\mathbb{R} \setminus \mathbb{Q}$ are equally dense in $\mathbb{R}$.
For any $x \in \mathbb{R}$ and $r > 0$, the interval $(x - r, x + r)$
is uncountable and its rational points are countable,
so it must contain irrationals.
Two disjoint sets can each be dense in the same space.
\end{remark*}

\begin{example}[Non-dense subset]
\label{ex:not-dense}
The integers $\mathbb{Z}$ are \emph{not} dense in $\mathbb{R}$.
The ball $B_{1/4}(1/2)$ is the interval $(1/4, 3/4)$, which contains
no integers.
Equivalently, $\overline{\mathbb{Z}} = \mathbb{Z} \neq \mathbb{R}$
(since $\mathbb{Z}$ is closed).
\end{example}

\begin{example}[Polynomials dense in $C([a,b])$]
\label{ex:poly-dense}
The Weierstrass Approximation Theorem states that $\mathbb{P}([a,b])$,
the set of polynomial functions on $[a,b]$, is dense in
$C([a,b])$ with the uniform metric $d(f,g) = \sup_{x \in [a,b]} |f(x) - g(x)|$.

This means: for every continuous $f$ and every $\varepsilon > 0$,
there is a polynomial $p$ with $d(f,p) < \varepsilon$.
Every continuous function can be uniformly approximated by a polynomial.
This is the precise statement behind the completion $\mathbb{P}([a,b]) \hookrightarrow C([a,b])$
mentioned in the real analysis bridge.
\end{example}

% ---- Definition: Separable ----

\begin{tcolorbox}[colback=propbox, colframe=propborder, arc=2pt,
  left=6pt, right=6pt, top=4pt, bottom=4pt,
  title={\small\textbf{Definition (Separable Metric Space)}},
  fonttitle=\small\bfseries]
\begin{definition}[Separable Metric Space]
\label{def:separable}
A metric space $(X, d)$ is \emph{separable} if it contains a
countable dense subset: there exists $D \subseteq X$ with
$D$ countable and $\overline{D} = X$.
\end{definition}
\end{tcolorbox}

\begin{remark*}[Fully quantified form]
\[
(X,d) \text{ separable}
\;\iff\;
\exists\, D \subseteq X : D \text{ is countable} \land \overline{D} = X.
\]
\end{remark*}

\begin{remark*}[Why ``separable''?]
The name is a historical accident and has nothing to do with
``separation axioms'' in topology.
Do not let the name mislead you.
Think of it instead as: a separable space can be
\emph{probed} or \emph{described} by a countable collection of points ---
every point in the space is a limit of a sequence from this countable subset.
\end{remark*}

\begin{remark*}[Why separability matters]
Separability is the correct notion of ``smallness'' or ``tractability''
for infinite-dimensional spaces.
Separable spaces are better behaved in analysis:
\begin{itemize}
  \item Every sequence has a subsequence whose indices we can name explicitly.
  \item Continuous functions on separable spaces are determined by their
        values on the dense subset.
  \item Many theorems (e.g.\ about compactness, integration, spectral theory)
        have cleaner statements or simpler proofs for separable spaces.
\end{itemize}
The spaces you encounter most in analysis --- $\mathbb{R}^n$,
$C([a,b])$, $\ell^p$ for $1 \le p < \infty$ --- are all separable.
Separability is what makes them ``manageable.''
\end{remark*}

\begin{proposition}[$\mathbb{R}^n$ is separable]
\label{prop:Rn-separable}
For each $n \ge 1$, the Euclidean space $(\mathbb{R}^n, d_2)$ is separable,
with dense subset $\mathbb{Q}^n = \{(q_1, \ldots, q_n) : q_i \in \mathbb{Q}\}$.
\end{proposition}

\begin{proof}
$\mathbb{Q}^n$ is countable (finite product of countable sets).
Let $x = (x_1, \ldots, x_n) \in \mathbb{R}^n$ and $\varepsilon > 0$.
For each coordinate $i$, choose $q_i \in \mathbb{Q}$ with
$|q_i - x_i| < \varepsilon / \sqrt{n}$.
Then $q = (q_1, \ldots, q_n) \in \mathbb{Q}^n$ satisfies
\[
d_2(q, x)
= \sqrt{\sum_{i=1}^n (q_i - x_i)^2}
< \sqrt{n \cdot \frac{\varepsilon^2}{n}}
= \varepsilon.
\]
So $B_\varepsilon(x) \cap \mathbb{Q}^n \neq \varnothing$ for every $x$ and $\varepsilon$,
giving $\overline{\mathbb{Q}^n} = \mathbb{R}^n$.
\end{proof}

\begin{proposition}[An uncountable discrete space is not separable]
\label{prop:discrete-not-sep}
Let $X$ be an uncountable set with the discrete metric.
Then $(X, d_{\mathrm{disc}})$ is not separable.
\end{proposition}

\begin{proof}
In the discrete metric, $B_{1/2}(x) = \{x\}$ for every $x \in X$.
If $D \subseteq X$ is any dense subset, then for each $x \in X$
the ball $B_{1/2}(x) = \{x\}$ must meet $D$, so $x \in D$.
Hence $D = X$, which is uncountable.
No countable dense subset exists.
\end{proof}

\begin{remark*}[Separability is a topological property]
Whether a metric space is separable depends only on its topology
(which sets are open), not on the specific distances.
Two equivalent metrics on the same set give the same open sets and
therefore agree on separability.
\end{remark*}

\begin{remark*}[Connection to the completion]
When a non-complete metric space $(X,d)$ is completed to $(\hat{X}, \hat{d})$,
the original space $X$ embeds as a dense subset of $\hat{X}$:
$\overline{X} = \hat{X}$.
This is exactly the density statement: every point in the completion
is a limit of a sequence in $X$.
The completion construction is the canonical way to ``fill in the gaps,''
and the resulting dense embedding is precisely what the co-Cauchy
equivalence classes represent.
\end{remark*}
